% ============================================================================
\clearpage
\section*{v4.8.0-r19 Source-Only: Frame Cubes, Transition Group and Holonomy}
\addcontentsline{toc}{section}{v4.8.0-r19 Source-Only: Frame Cubes, Transition Group and Holonomy}
\label{sec:v480-r19-frame-cubes-holonomy}
% ============================================================================

\paragraph{Purpose.}
Layer r18 typed the frame offset by cross terms between the three frames.  Two working results followed outside the
source: the pre-registered test H12 (\texttt{llm/Ergebnis\_H12\_Drei\_Frame\_Zuordnung.md}) and a backward analysis
with a readout index (\texttt{llm/k34\_rueckwaerts\_3x3x3.py}, map entry K39).  This layer records both, asks what a
family of readout cubes linked by transition maps adds, and compares the transition group of the earlier reading with
that of Reading~B (\S\ref{subsec:v480-r06-readings}).  It contains typing and arithmetic, no new prediction.  Earlier
layers are not rewritten.

\noindent\textbf{Source-only status.}
This v4.8.0-r19 layer is source-only.  It does not build a PDF, does not assign a new DOI,
and does not authorize a public release.  The v4.8.0 Version DOI
\href{https://doi.org/10.5281/zenodo.23365244}{10.5281/zenodo.23365244} and the Concept DOI
\href{https://doi.org/10.5281/zenodo.19210728}{10.5281/zenodo.19210728} are unchanged by it.

\subsection*{v4.8.0-r19 records of H12 and of the backward analysis}
\label{subsec:v480-r19-records}

\paragraph{H12.}
With \(\cosh(\eta_++\eta_s)=\gL\), \(\Omega_M=\tanh^2\eta_s\) (Planck 2018 with BAO) and
\(\sin^2\theta_{12}=\tanh^2\eta_+/\cos^2\theta_{13}\), the assignment gives \(\beta_+=0.5408\), \(\beta_s=0.5578\) and
\(\sin^2\theta_{12}=0.2991\pm0.0057\), against \(0.3086\) for the symmetric relation H9.  At the precision fixed for
the next JUNO release the two are \(1.36\sigma\) apart; by the fixed criterion they cannot be separated, and the
result is recorded as no test.  The limit is set by the uncertainty of \(\Omega_M\), not by JUNO.

\paragraph{Backward analysis (does not count as a test).}
Let \(T_{abc}\) be the offset of frame \(a\) against frame \(b\), read from frame \(c\), starting from the Hopf value
\(\bsym=1/\sqrt3\).  Pure kinematics makes \(T_{abc}=\eta_a-\eta_b\) independent of \(c\), since a boost into \(c\)
shifts all rapidities equally.  Of the readout terms tried, those depending on the distance between reading and read
frame fail, because they vanish for a self-reading (\(\chi^2\approx34\)).  A single factor on everything read from
the matter frame, \(T_{abc}=(\eta_a-\eta_b)f(c)\), fits the two measurements with \(f(+)=0.944\pm0.008\)
(\(\chi^2=2.9\), one degree of freedom).  In \(\beta\) this is \(0.956\pm0.007\), numerically the known Hopf gap
\(\delta_{\mathrm{CL}}=0.951\) of Remark~\ref{conj:cl-correction}.  The backward analysis yields no new number; it
offers a different reading of a known one.  Only \(f(+)\) is accessible, since every measurement is made from the
matter frame.

\subsection*{v4.8.0-r19 cubes linked by transition maps}
\label{subsec:v480-r19-cubes}

Take one readout cube \(T^{(d)}_{abc}\) per frame \(d\in\{+,0,-\}\) and transition maps
\(\Lambda(d'\!\leftarrow\!d)\) with \(T^{(d')}=\Lambda(d'\!\leftarrow\!d)\,T^{(d)}\).  If the maps are fixed by the
theory, the three cubes carry the information of one cube and the transition law; the family adds no parameter.
What it adds is a consistency question: whether the maps compose to the identity around a closed path of frames.  The
residue, if any, is the holonomy of the family.  It depends on the transition group and the positions of the frames
only, and it is the part that no choice of the reading frame can remove.  If the transition maps commute, the
holonomy is trivial and the family adds nothing.

\subsection*{v4.8.0-r19 the transition group in the two readings}
\label{subsec:v480-r19-transition-group}

\paragraph{Earlier reading.}
Theorem~\ref{thm:symmetric-frame} places matter and tachyon at \(\pm\bsym\) on one axis around the photon and calls
\(\gL\) the Lorentz factor between them.  The transitions are collinear boosts; they commute, and the holonomy is
trivial.

\paragraph{Reading~B.}
By \S\ref{subsec:v480-r06-readings} the kernel is the Riemann--Silberstein vector \(z=E+iB\) under
\(SO(3,\mathbb C)\cong SO^+(1,3)\), with \(s=E^2-B^2+2i\,E\cdot B\).  The working assignment of r07 reads matter as
\(s>0\) (electric type), tachyon as \(s<0\) (magnetic type), photon as \(s=0\) (null field).

\paragraph{Theoremlet v4.8.0-r19.1.}
In Reading~B no Lorentz transformation relates matter to tachyon, and none relates either of them to the photon.
Matter and tachyon with equal \(|s|\) are related by the duality rotation \(z\mapsto e^{i\pi/2}z\), that is
\(E\mapsto-B\), \(B\mapsto E\).  The photon is reached by neither group action.

\paragraph{Proof sketch.}
\(SO(3,\mathbb C)\) preserves \(s\) (Theoremlet r06.3), so it preserves its sign and the null cone.  The duality
rotation multiplies \(s\) by \(e^{2i\alpha}\) and maps \(s>0\) to \(s<0\) at \(\alpha=\pi/2\)
(Theoremlet r07.2); it preserves \(s=0\).  This is the familiar statement that a radiation field is null in every
frame and under every duality rotation.  \(\square\)

\paragraph{Consequence.}
The Lorentz factor \(\gL\) between matter and tachyon in Theorem~\ref{thm:symmetric-frame} has no counterpart in
Reading~B, where the two are related by a duality rotation and not by a boost.  On this point the two readings
disagree about the transition group itself.

\paragraph{Theoremlet v4.8.0-r19.2.}
The duality circle carries a \(\mathbb Z_2\) holonomy around the photon.  Under \(z\mapsto e^{i\alpha}z\) one has
\(s\mapsto e^{2i\alpha}s\); as \(\alpha\) runs from \(0\) to \(\pi\), \(s\) encircles \(s=0\) once and \(z\) returns as
\(-z\).  Paths of \(s\) that do not encircle \(s=0\) have trivial holonomy.  The photon is the branch point of
\(z\) over \(s\).

\paragraph{Proof sketch.}
\(s\) is quadratic in \(z\), so the map from the duality circle of \(z\) to the circle of \(s\) is two to one.
Theoremlet r07.1 states the sign ambiguity of \(z_m\) under \(\alpha\mapsto\alpha+\pi\).  \(\square\)

\subsection*{v4.8.0-r19 the earlier reading as a photon-centred view}
\label{subsec:v480-r19-photon-centred}

\paragraph{Theoremlet v4.8.0-r19.3.}
The real axis of the \(s\)-plane reproduces the sector structure of Theorem~\ref{thm:symmetric-frame}: matter and
tachyon at \(\pm|s|\), mirror-symmetric about the photon at \(s=0\), the mirror \(s\mapsto-s\) being the half-turn of
the duality circle.

\paragraph{Proof sketch.}
Theoremlet r07.1 with \(\alpha=0\) and \(\alpha=\pi/2\).  \(\square\)

\paragraph{Reading, offered and not adopted.}
The photon is fixed by the whole duality circle, so it picks out no reference direction in the \(s\)-plane.  The
earlier reading can then be taken as the view centred on the photon with the real axis chosen as reference, that is,
with the duality angle fixed as a gauge.  On this axis the straight path from matter to tachyon passes through the
photon, the branch point, where the phase is undefined; in Reading~B the path runs around the photon on
\(|s|=\mathrm{const}\).  The photon-centred view therefore cannot see the \(\mathbb Z_2\) holonomy.  It also supplies
something Reading~B lacks: a metric along the axis, in which the mirror distance is read as a boost with factor
\(\gL\).  No correspondence between \(|s|\) and \(\bsym\) is given in the source, and by r19.1 the boost reading has no
counterpart in Reading~B.  Whether the duality angle is a gauge, which would make the earlier reading a gauge-fixed
projection, or physical, which would make it a projection that loses information, is the question left open in r07.

\subsection*{v4.8.0-r19 consequence for the readout factor}
\label{subsec:v480-r19-readout-factor}

The \(\mathbb Z_2\) holonomy can only contribute \(\pm1\); it cannot produce a factor such as \(f(+)=0.944\).  A
continuous holonomy would need closed paths in an \(SO(3,\mathbb C)\)-orbit together with a rule for transporting
frames along them, i.e.\ a connection.  The source supplies neither.  Composing boosts in non-collinear directions
yields a Wigner rotation, and the triole directions are a candidate for such directions, but choosing them is a model
decision.  Any comparison of such a holonomy with \(f(+)\) is a prediction and requires a pre-registration note.

\subsection*{v4.8.0-r19 register entries}
\label{subsec:v480-r19-register}

\begin{center}
\small
\begin{tabular}{@{}>{\raggedright\arraybackslash}p{0.34\textwidth}>{\raggedright\arraybackslash}p{0.58\textwidth}@{}}
\toprule
statement & status after r19\\
\midrule
Thm.~\ref{thm:symmetric-frame} & sector structure reproduced by the real \(s\)-axis of Reading~B (r19.3); \(\gL\) as a Lorentz factor between matter and tachyon has no counterpart in Reading~B (r19.1)\\
duality circle of r07 & \(\mathbb Z_2\) holonomy around the photon, photon as branch point (r19.2)\\
Rem.~\ref{conj:cl-correction}, \(\delta_{\mathrm{CL}}\) & numerically equal to the readout factor \(f(+)\) of the backward analysis; reinterpretation only, not a test (r19)\\
H12 & not separable from H9 by the fixed criterion; no test (r19)\\
\bottomrule
\end{tabular}
\end{center}

\subsection*{v4.8.0-r19 ledger and guard update}
\label{subsec:v480-r19-ledger-guard}

\[
R_{\rm hol}^{v480\text{-}r19}=(0,0,1)
\]
for (transition groups of both readings compared, \(\mathbb Z_2\) holonomy of the duality circle recorded,
continuous holonomy open for lack of a connection).

The layer spends no guard:
\[
R_{\rm glob}^{(v480\text{-}r19,+)}=(0,1,1,1,1),
\qquad
\Gamma_{\rm guard}^{v480\text{-}r19}=(1,1,1,1,1),
\qquad
\beta_4=1/256\ \text{unspent}.
\]
It does not imply exact finite \(C5\) zero closure, exact all-readout terminality,
\(C6/T2\), a global photon-clock/RH theorem, or global all-readout closure.

